\section{Ôn tập giữa kì 2}
\subsection{Đề 1}
\Opensolutionfile{ans}[ans/ans-2-OTGK2-De1]
\caulc
\begin{ex}%[2D4N1-1]
	Hàm số $F(x)$ là một nguyên hàm của hàm số $f(x)$ trên khoảng $\mathbb{K}$ nếu
	\choice
	{$F'(x)=-f(x),\forall x\in \mathbb{K}$}
	{$f'(x)=F(x),\forall x\in \mathbb{K}$}
	{\True $F'(x)=f(x),\forall x\in \mathbb{K}$}
	{$f'(x)=-F(x),\forall x\in \mathbb{K}$}
	\loigiai{
		Hàm số $F(x)$ là một nguyên hàm của hàm số $f(x)$ trên khoảng $\mathbb{K}$ nếu $F'(x)=f(x),\forall x\in \mathbb{K}$.
	}
\end{ex}
%Câu 2
\begin{ex}%[2D4H1-2]
	Tìm một nguyên hàm $F(x)$ của $f(x)=x-x^{-2},$ biết $F(1)=0$.
	\choice
	{$F(x)=\dfrac{x^{2}}{2}+\dfrac{1}{x}+\dfrac{3}{2}$}
	{$F(x)=\dfrac{x^{2}}{2}-\dfrac{1}{x}+\dfrac{1}{2}$}
	{\True $F(x)=\dfrac{x^{2}}{2}+\dfrac{1}{x}-\dfrac{3}{2}$}
	{$F(x)=\dfrac{x^{2}}{2}-\dfrac{1}{x}-\dfrac{1}{2}$}
	\loigiai{
		Ta có $\displaystyle  \int (x-x^{-2})\ \mathrm{d}x = \dfrac{x^2}{2} + \dfrac{1}{x} +C$.  \\
		Mà $F(1) =0 \Leftrightarrow \dfrac{1}{2} +1 +C = 0 \Leftrightarrow  C= -\dfrac{3}{2}$. \\
		Do đó $F(x) = \dfrac{x^{2}}{2}+\dfrac{1}{x}-\dfrac{3}{2}$.
	}
\end{ex}
%Câu 3
\begin{ex}%[2D4N2-1]
	Xét $ f(x)$ là một hàm số tùy ý, $ F(x)$ là một nguyên hàm của $f(x)$ trên đoạn $\left[a;\,b\right]$. Mệnh đề nào dưới đây đúng?
	\choice
	{\True $\displaystyle\int\limits_a^b{f(x)\mathrm{\,d}x}=F(b)-F(a)$}
	{$\displaystyle\int\limits_a^b{f(x)\mathrm{\,d}x}=F(a)-F(b)$}
	{$\displaystyle\int\limits_a^b{f(x)\mathrm{\,d}x}=F(a)+F(b)$}
	{$\displaystyle\int\limits_a^b{f(x)\mathrm{\,d}x}=-F(a)-F(b)$}
	\loigiai{
		Mệnh đề đúng là $\displaystyle\int\limits_a^b{f(x)\mathrm{\,d}x}=F(b)-F(a)$.
	}
\end{ex}
%Câu 4
\begin{ex}%[2D4N1-1]
	Nếu $\displaystyle\int\limits_{2}^{5}f(x)\mathrm{\,d}x=4$ và $\displaystyle\int\limits_{2}^{5}g(x)\mathrm{\,d}x=5$ thì $\displaystyle\int\limits_{2}^{5}[2 f(x)+g(x)]\mathrm{\,d}x$ bằng
	\choice
	{\True $13$}
	{$3$}
	{$-1$}
	{$-3$}
	\loigiai{Ta có $\displaystyle\int\limits_{2}^{5}[2 f(x)+g(x)]\mathrm{\,d}x=2\displaystyle\int\limits_{2}^{5}f(x)\mathrm{\,d}x+\displaystyle\int\limits_{2}^{5}g(x)\mathrm{\,d}x=2\cdot 4+5=13$.}
\end{ex}
%Câu 5
\begin{ex}%[2D4N3-1]
	Gọi $S$ là diện tích hình phẳng giới hạn bởi các đường $y=3^x$, $y=0$, $x=0$ và $x=2$. Mệnh đề nào dưới đây là đúng?
	\choice
	{$S=\pi \displaystyle\int\limits_{0}^{2}{3^x}\mathrm{\,d}x$}
	{$S=\displaystyle\int\limits_{0}^{2}{3^{2x}}\mathrm{\,d}x$}
	{\True $S=\displaystyle\int\limits_{0}^{2}{3^x}\mathrm{\,d}x$}
	{$S=\pi \displaystyle\int\limits_{0}^{2}{3^{2x}}\mathrm{\,d}x$}
	\loigiai
	{Diện tích hình phẳng giới hạn bởi các đường $y=3^x$, $y=0$, $x=0$ và $x=2$ là $S=\displaystyle\int\limits_{0}^{2}{\left|3^x\right |}\mathrm{\,d}x = \displaystyle\int\limits_{0}^{2}{3^x}\mathrm{\,d}x$.}
\end{ex}
%Câu 6
\begin{ex}%[2D4H3-3]
	Kí hiệu $(H)$ là hình phẳng giới hạn bởi đồ thị hàm số $y=2x-x^2$ và $y=0$. Tính thể tích vật thể tròn xoay được sinh ra bởi hình phẳng $(H)$ khi nó quay quanh trục $Ox$.
	\choice
	{$\dfrac{17\pi}{15}$}
	{$\dfrac{19\pi}{15}$}
	{\True $\dfrac{16\pi}{15}$}
	{$\dfrac{18\pi}{15}$}
	\loigiai{Phương trình hoành độ giao điểm của hai đồ thị hàm số là $2x-x^2=0\Leftrightarrow\hoac{&x=0\\&x=2.}$\\
		Thể tích khối tròn xoay sinh ra là
		\[V=\pi\displaystyle\int\limits_0^2\left(2x-x^2\right)^2\mathrm{\,d}x=\pi\displaystyle\int\limits_0^2\left(x^4-4x^3+4x^2\right)\mathrm{\,d}x=\pi\left(\dfrac{x^5}{5}-x^4+\dfrac{4x^3}{3}\right)\Bigg|_0^2=\dfrac{16\pi}{15}.\]}
\end{ex}
%Câu 7
\begin{ex}%[2H5N1-2]
	Trong không gian $Oxyz$, mặt phẳng $(P)\colon x + y + z +1=0$ có một véc-tơ pháp tuyến là
	\choice
	{$\overrightarrow{n}_1=(-1;1;1)$}
	{$\overrightarrow{n}_4=(1;1;-1)$}
	{\True $\overrightarrow{n}_3=(1;1;1)$}
	{$\overrightarrow{n}_2=(1;-1;1)$}
	\loigiai{
		$(P)\colon x + y + z +1=0$ có một véc-tơ pháp tuyến là	$\overrightarrow{n}_3=(1;1;1)$.
	}
\end{ex}
%Câu 8
\begin{ex}%[2H5N1-3]
	Trong không gian $Oxyz$, cho ba điểm $A(2;0;0)$, $B(0;-1;0)$, $C(0;0;3)$. Mặt phẳng $(ABC)$ có phương trình là
	\choice
	{$\dfrac{x}{2}+\dfrac{y}{1}+\dfrac{z}{3}=1$}
	{$\dfrac{x}{-2}+\dfrac{y}{1}+\dfrac{z}{3}=1$}
	{$\dfrac{x}{2}+\dfrac{y}{1}+\dfrac{z}{-3}=1$}
	{\True $\dfrac{x}{2}+\dfrac{y}{-1}+\dfrac{z}{3}=1$}
	\loigiai{
		Ta thấy ba điểm $A$, $B$, $C$ lần lượt thuộc các trục tọa độ $Ox$, $Oy$, $Oz$ nên mặt phẳng $(ABC)$ có phương trình mặt phẳng theo đoạn chắn là
		\[\dfrac{x}{2}+\dfrac{y}{-1}+\dfrac{z}{3}=1.\]
	}
\end{ex}
%Câu 9
\begin{ex}%[2H5H1-4]
	Trong không gian $Oxyz$, cho hai mặt phẳng $(P)\colon x-2y-z+2=0$ và $(Q)\colon 2x-3y+mz+1=0$ ( $m$ là tham số). Tìm $m$ để $(P) \perp (Q)$.
	\choice
	{\True $m=8$}
	{$m=-4$}
	{$m=-8$}
	{$m=4$}
	\loigiai{
		Mặt phẳng $(P)\colon x-2y-z+2=0$ có một véc-tơ pháp tuyến là $\overrightarrow{n}_P=(1;-2;-1)$.\\
		Mặt phẳng $(Q)\colon 2x-3y+mz+1=0$ có một véc-tơ pháp tuyến là $\overrightarrow{n}_Q=(2;-3;m)$.\\
		Do đó, $(P) \perp (Q)$ khi và chỉ khi
		\[ \overrightarrow{n}_P\cdot \overrightarrow{n}_Q=0 \Leftrightarrow 1\cdot 2+(-2)\cdot (-3)+(-1)\cdot m=0 \Leftrightarrow 8-m=0 \Leftrightarrow m=8. \]
	}
\end{ex}
%Câu 10
\begin{ex}%[2H5N2-1]
	Trong không gian $Oxyz$, điểm nào dưới đây thuộc đường thẳng $d\colon\heva{& x=1+2t\\ & y=-3+t\\ & z=4+5t}$?
	\choice
	{$Q(4;1;3)$}
	{$P(3;-2;-1)$}
	{$N(2;1;5)$}
	{\True $M(1;-3;4)$}
	\loigiai{
		Dễ thấy đường thẳng $d$ đi qua điểm $M(1;-3;4)$.
	}
\end{ex}
%Câu 11
\begin{ex}%[2H5N2-3]
	Trong không gian $Oxyz$, đường thẳng đi qua điểm $I(1;-1;-1)$ và nhận $\overrightarrow{u}=(-2;3;-5)$ làm véc-tơ chỉ phương có phương trình chính tắc là
	\choice
	{\True $\dfrac{x-1}{-2} = \dfrac{y+1}{3} = \dfrac{z+1}{-5}$}
	{$\dfrac{x-1}{-2} = \dfrac{y+1}{3} = \dfrac{z+1}{5}$}
	{$\dfrac{x+2}{1} = \dfrac{y-3}{-1} = \dfrac{z+5}{-1}$}
	{$\dfrac{x+1}{-2} = \dfrac{y-1}{3} = \dfrac{z-1}{-5}$}
	\loigiai
	{Đường thẳng đi qua điểm $I(1;-1;-1)$ và nhận $\overrightarrow{u}=(-2;3;-5)$ làm véc-tơ chỉ phương có phương trình chính tắc là $\dfrac{x-1}{-2} = \dfrac{y+1}{3} = \dfrac{z+1}{-5}$.}
\end{ex}
%Câu 12
\begin{ex}%[2H5H2-7]
	Trong không gian $Oxyz$, tính góc giữa hai đường thẳng $d_1\colon \dfrac{x}{1}=\dfrac{y+1}{-1}=\dfrac{z-1}{2}$ và $d_2\colon\dfrac{x+1}{-1}=\dfrac{y}{1}=\dfrac{z-3}{1}$.
	\choice
	{$30^{\circ}$}
	{\True $90^{\circ}$}
	{$45^{\circ}$}
	{$60^{\circ}$}
	\loigiai
	{
		$d_1	$ có một véc-tơ chỉ phương $\overrightarrow{u_1}=(1;-1;2)$; $d_2$ có một véc-tơ chỉ phương $\overrightarrow{u_2}=(-1;1;1)$.\\
		Vì $\overrightarrow{u_1}\cdot\overrightarrow{u_2}= 1\cdot (-1)+(-1)\cdot 1+2\cdot 1=0$ nên $\overrightarrow{u_1}\perp\overrightarrow{u_2}$.\\
		Vậy góc giữa hai đường thẳng $d_1$ và $d_2$ bằng $90^{\circ}$.
	}
\end{ex}
\Closesolutionfile{ans}
\indapan{6}{ans/ans-2-OTGK2-De1}

\cauds
\Opensolutionfile{ans}[ans/ans-2-OTGK2-De1-DS]
\begin{ex}%[2025-TLOT-2018,Trần Xuân Hòa]%[2D4H2-2]
	Cho hàm số $f(x)=x(x^2+3)$. Xét $I=\displaystyle\int\limits_{-1}^1|f(x)|\mathrm{\; d}x$.
	\choiceTF
	{\True Đặt $I_1=\displaystyle\int\limits_{0}^1|f(x)|\mathrm{\; d}x$ và $I_2=\displaystyle\int\limits_{-1}^0|f(x)|\mathrm{\; d}x$. Khi đó $I_1=I_2$}
	{Giá trị $I=0$}
	{\True Số thực dương $m$ để $\displaystyle\int\limits_0^m|f(x)|\mathrm{\;d}x=4$ bằng $\sqrt{2}$}
	{Số thực $a$ để $\displaystyle\int\limits_0^1x(x^2+3-a\sqrt{x})\mathrm{\; d}x=0$ bằng $4$}
	\loigiai{
	\begin{itemchoice}
		\itemch Đúng. 
		\begin{itemize}
			\item Ta có $I_1=\displaystyle\int\limits_0^1(x^3+3x)\mathrm{\; d}x=\left(\dfrac{x^4}{4}+\dfrac{3x^2}{2}\right)\bigg |_0^1=\dfrac{7}{4}$.
			\item Ta có $I_2=\displaystyle\int\limits_{-1}^0(-x^3-3x)\mathrm{\; d}x=\left(-\dfrac{x^4}{4}-\dfrac{3x^2}{2}\right)\bigg |_{-1}^0=\dfrac{7}{4}$.
		\end{itemize}
		Do đó $I_1=I_2$.
		\itemch Sai. Ta có $I=\displaystyle\int\limits_{-1}^0|f(x)|\mathrm{\; d}x+\displaystyle\int\limits_{0}^1|f(x)|\mathrm{\; d}x=I_1+I_2=\dfrac{7}{4}+\dfrac{7}{4}=\dfrac{7}{2}$.
		\itemch Đúng. Ta có
		\begin{eqnarray*}
			&&\displaystyle\int\limits_0^m|f(x)|\mathrm{\;d}x=4\\
			&\Leftrightarrow&\displaystyle\int\limits_0^{m}(x^3+3x)\mathrm{\;d}x=4\\
			&\Leftrightarrow&\left(\dfrac{x^4}{4}+\dfrac{3x^2}{2}\right)\bigg |_0^m=4\\
			&\Leftrightarrow&\dfrac{m^4}{4}+\dfrac{3m^2}{2}-4=0\\
			&\Leftrightarrow&\hoac{&m=\sqrt{2}\\&m=-\sqrt{2} \text{ ( loại )}.}
		\end{eqnarray*}
		Vậy $m=\sqrt{2}$ thỏa mãn.
		 \itemch Sai. Ta có
		\begin{eqnarray*}
			&&\displaystyle\int\limits_0^1x(x^2+3-a\sqrt{x})\mathrm{\; d}x=0\\
			&\Leftrightarrow&\displaystyle\int\limits_0^1\left(x^3+3x-ax^{\tfrac{3}{2}}\right)\mathrm{\; d}x=0\\
			&\Leftrightarrow&\left(\dfrac{1}{4}x^4+\dfrac{3}{2}x^2-\dfrac{2a}{5}x^{\tfrac{5}{2}}\right)\bigg|_0^1=0\\
			&\Leftrightarrow&\dfrac{7}{4}-\dfrac{2a}{5}=0\\
			&\Leftrightarrow&a=\dfrac{35}{8}.
		\end{eqnarray*}
	\end{itemchoice}
	}
\end{ex}
\begin{ex}%[2025-TLOT-2018,Trần Xuân Hòa]%[2D4V3-3]
	\immini{Gọi $(H)$ là hình giới hạn bởi đồ thị các hàm số $y=\sqrt{x}$, $y=2-x$ và trục hoành. Kí hiệu diện tích hình $(H)$ là $S_1$ và diện tích hình phẳng giới hạn bởi đồ thị hàm số $y=2-x$, $y=\sqrt{x}$ và trục $Oy$ là $S_2$.
	}{\begin{tikzpicture}[scale=0.8, font=\footnotesize, line join = round, line cap = round, >=stealth]
		\draw[->](-1,0) -- (0,0) node[below left]{$O$} -- (4,0) node[below]{$x$};
		\draw[->](0,-1.5) -- (0,3) node[right]{$y$};
		\foreach \x in {1,2} \draw (\x,0) node[below]{$\x$} circle(1pt);	
		\foreach \y in {1,2} \draw (0,\y) node[left]{$\y$} circle(1pt);	
		\draw[domain=0:3,samples=100] plot(\x,{(\x)^0.5}) node[below]{$y=\sqrt{x}$};	
		\draw[domain=-0.5:3,samples=100] plot(\x,{2-\x}) node[left]{$y=2-x$};		
		\draw[pattern = north east lines, draw=none] (0,0) plot[domain=0:1] (\x, {(\x)^0.5})--(2,0)--cycle;
		\draw[dashed] (1,1)--(1,0);
		\end{tikzpicture}}
	\choiceTF
	{\True Thể tích khối tròn xoay khi quay hình phẳng giới hạn bởi đồ thị hàm số $y=2-x$, $x=0$, $x=1$ và trục $Ox$ xung quanh trục $Ox$ bằng $\dfrac{7\pi}{3}$}
	{\True Giá trị $S_1=\dfrac{7}{6}$}
	{$S_1=S_2$}
	{Thể tích khối tròn xoay được tạo bởi khi quay hình $(H)$ quanh trục $Ox$ bằng $\pi$}
	\loigiai{
		
\begin{itemchoice}
	\itemch Đúng. Ta có $V=\pi\displaystyle\int\limits_0^1(2-x)^2\mathrm{\;d}x=\dfrac{7\pi}{3}$.
	\itemch Đúng. 	Phương trình hoành độ giao điểm \\$\sqrt{x}=2-x \Leftrightarrow \left\{\begin{aligned}
		&0\leqslant x\leqslant 2 \\
		&x=4-4x+x^2
	\end{aligned}\right. \Leftrightarrow x=1$. \\
	Vậy $S=\displaystyle\int\limits_0^1 \sqrt{x} \mathrm{\,d}x+\displaystyle\int\limits_1^2 (2-x) \mathrm{\,d}x=\dfrac{2}{3}x\sqrt{x}\bigg|_0^1+\left(2x-\dfrac{x^2}{2}\right)\bigg|_1^2=\dfrac{2}{3}+\dfrac{1}{2}=\dfrac{7}{6}$.
	\itemch Sai. Giả sử đồ thị hàm số $y=2-x$ cắt $Ox$, $Oy$ lần lượt tại $A$, $B$.\\
	 Diện tích tam giác $OAB$ bằng $S=\dfrac{1}{2}OA\cdot OB=2$. \\
	Giá trị $S_2=S-S_1=2-\dfrac{7}{6}=\dfrac{5}{6}$. Do đó $S_1>S_2$.
	\itemch Sai. Xét các phương trình sau
	\begin{itemize}
		\item $\sqrt{x} = 0 \Leftrightarrow x=0$.
		\item $\sqrt{x} = 2-x \Leftrightarrow \left(\sqrt{x}\right)^2 + \sqrt{x} -2 =0 \Leftrightarrow x=1$.		
		\item $2-x=0 \Leftrightarrow x=2$.
	\end{itemize}
	Thể tích của khối tròn xoay cần tính là 	
	\allowdisplaybreaks
	\begin{eqnarray*}
		V &=& \pi \displaystyle \int\limits_{0}^{1} \left(\sqrt{x}\right)^2 \mathrm{\,d}x + \pi \displaystyle \int\limits_{1}^{2} (2-x)^2 \mathrm{\,d}x = \pi \displaystyle \int\limits_{0}^{1} x \mathrm{\,d}x + \pi \displaystyle \int\limits_{1}^{2} \left(4-4x+x^2\right) \mathrm{\,d}x\\		
		&=& \pi	\cdot \dfrac{x^2}{2} \bigg|_{0}^{1}	 + \pi	\cdot \left(4x - 2x^2 + \dfrac{x^3}{3}\right) \bigg|_{1}^{2} = \dfrac{\pi}{2} + \dfrac{\pi}{3} = \dfrac{5\pi}{6}.
	\end{eqnarray*}
\end{itemchoice}
}
\end{ex}

\begin{ex}%[2025-TLOT-2018,Trần Xuân Hòa]%[2H5H1-3]
	Trong không gian với hệ tọa độ $Oxyz$, cho các điểm $A(2;0;0)$, $B(0;3;0)$, $C(0;0;6)$, $D(1;1;1)$.
	\choiceTF
	{Phương trình mặt phẳng trung trực của đoạn thẳng $AC$ là $x-3z-8=0$}
	{\True Phương trình mặt phẳng đi qua $D$ và chứa trục $Oz$ là $x-y=0$}
	{\True Tọa độ điểm $M$ có hoành độ dương thuộc $Ox$ sao cho khoảng cách từ $M$ đến mặt phẳng $(ACD)$ bằng $\sqrt{14}$ là $M\left(\dfrac{20}{3};0;0\right)$}
	{Có tất cả $5$ mặt phẳng phân biệt đi qua $3$ trong $5$ điểm $O$,$A$,$B$,$C$,$D$}
	\loigiai{
		\begin{itemchoice}
			\itemch Sai. Véc-tơ $\overrightarrow{AC}=(-2;0;6)$. Gọi $I$ là trung điểm của $AC$. Khi đó $I(1;0;3)$. \\
			Mặt phẳng trung trực của $AC$ đi qua $I$ và nhận véc-tơ $\overrightarrow{IA}=(1;0;-3)$ làm véc-tơ pháp tuyến có phương trình $1\cdot (x-1)+0\cdot (y-0)+(-3)\cdot (z-3)=0$ hay $x-3z+8=0$.\\
			Vậy phương trình mặt phẳng trung trực của đoạn thẳng $AC$ là $x-3z+8=0$.	
			\itemch Đúng. Véc-tơ $\overrightarrow{OD}=(1;1;1)$. Gọi $E(0;0;1)\in Oz$, $\overrightarrow{OE}=(0;0;1)$ .\\
			 Mặt phẳng $(Q)$ cần tìm có cặp véc-tơ chỉ phương là $\overrightarrow{OD}$ và $\overrightarrow{OE}$ nên có véc-tơ pháp tuyến $\left[\overrightarrow{OD},\overrightarrow{OE}\right]=(1;-1;0)$.\\
			Mặt phẳng $(Q)$ đi qua $D$ và có véc-tơ pháp tuyến $(1; -1; 0)$ nên có phương trình $1\cdot (x-1)-1\cdot (y-1)+0\cdot (z-1)=0$ hay $x-y=0$.\\
			Vậy phương trình mặt phẳng đi qua $D$ và chứa trục $Oz$ là $x-y=0$.
			\itemch Đúng. Ta thấy $3$ điểm $A,B,C$ tạo thành mặt phẳng đoạn chắn các trục tọa độ có phương trình 	$\dfrac{x}{2}+\dfrac{y}{3}+\dfrac{z}{6}=1$ hay $3x+2y+z-6=0$.\\
			Suy ra $D\in (ABC)$. Như vậy $4$ điểm $A,B,C,D$ đồng phẳng. \\
			Gọi $M(m;0;0)\in Ox$. Ta có 
			\begin{eqnarray*}
				&&\mathrm{d}(M,ACD)=\sqrt{14}\\
				&\Leftrightarrow&\dfrac{|3\cdot m+2\cdot 0+0-6|}{\sqrt{3^2+2^2+1^2}}=\sqrt{14}\\
				&\Leftrightarrow&|3m-6|=14\\
				&\Leftrightarrow&\hoac{&m=\dfrac{20}{3}\\&m=-\dfrac{8}{3} \text{ (loại)}.}
			\end{eqnarray*}
			Vậy $M\left(\dfrac{20}{3};0;0\right)$.
			\itemch Sai. Ta có $A$, $B$, $C$, $D$ đồng phẳng. Qua $3$ điểm phân biệt không thẳng hàng ta xác định được $1$ mặt phẳng.
			\begin{itemize}
				\item Số mặt phẳng đi qua $3$ trong $5$ điểm trên là $\mathrm{C}_5^3$.
				\item Bốn mặt phẳng trùng nhau là $(ABC)$, $(ABD)$, $(BCD)$, $(ACD)$.
			\end{itemize}
			Vậy nên số mặt phẳng phân biệt đi qua $3$ trong $5$ điểm $O,A,B,C,D$ là $\mathrm{C}_5^3-4+1=7$.\\
			Vậy có $7$ mặt phẳng phân biệt đi qua $3$ trong $5$ điểm $O$,$A$,$B$,$C$,$D$.
		\end{itemchoice}
	}
\end{ex}
\begin{ex}%[2025-TLOT-2018,Trần Xuân Hòa]%[2H5V2-3]
		Trong không gian $Oxyz$, cho mặt phẳng $(P)\colon x+2y-3z-12=0$ và đường thẳng $d$ có phương trình $d\colon  \dfrac{x+7}{3}=\dfrac{y+10}{4}=\dfrac{z-4}{-2}$.
	\choiceTF
	{\True Toạ độ giao điểm $M$ của đường thẳng $d$ với mặt phẳng $(P)$ là $M(2;2;-2)$}
	{\True Gọi $\varphi $ là góc giữa $d$ và $(P)$. Khi đó $\sin \varphi=\dfrac{17\sqrt{406}}{406}$}
	{Phương trình đường thẳng $d'$ là hình chiếu vuông góc của $d$ lên mặt phẳng $\left(Oxy\right)$ là $\heva{&x=-1+3t\\&y=2+4t\\&z=0}$}
	{Đường thẳng nằm trong mặt phẳng $(P)$, đồng thời cắt và vuông góc với đường thẳng $d$ có phương trình là $\dfrac{x-2}{8}=\dfrac{y-2}{-7}=\dfrac{z-2}{2}$}
	\loigiai{Đặt $\dfrac{x+7}{3}=\dfrac{y+10}{4}=\dfrac{z-4}{-2}=t$ nên đường thẳng $d$ có phương trình tham số $d\colon \heva{&x=-7+3t\\&y=-10+4t\\&z=4-2t}$.
		\begin{itemchoice}
			\itemch Đúng. Toạ độ giao điểm của $d$ và $(P)$ là nghiệm của hệ phương trình: \\
			$$\heva{&x=-7+3t\qquad (1)\\&y=-10+4t \qquad (2)\\&z=4-2t \qquad (3)\\&x+2y-3z-12=0\qquad (4)}$$
			Thay $(1)$, $(2)$, $(3)$ vào $(4)$ ta được $t=3$. \\
			Vậy $M(2;2;-2)$ là giao điểm của đường thẳng $d$ với mặt phẳng $(P)$.
			\itemch Đúng. Một véc-tơ chỉ phương của $d$ là $\overrightarrow{u}=(3; 4; -2)$, một véc-tơ pháp tuyến của $(P)$ là $\overrightarrow{n}=(1; 2; -3)$.\\
			Ta có $\sin\varphi=\dfrac{|\overrightarrow{n}\cdot \overrightarrow{u}|}{|\overrightarrow{n}|\cdot |\overrightarrow{u}|}=\dfrac{|3\cdot 1+4\cdot 2+(-2)\cdot (-3)|}{\sqrt{3^2+4^2+(-2)^2}\cdot \sqrt{1^2+2^2+(-3)^2}}=\dfrac{17\sqrt{406}}{406}$.
			\itemch Sai. Gọi $N$ là giao điểm của $d$ và mặt phẳng $(Oxy)\colon z=0$.\\
			 Tọa độ $N$ là nghiệm của hệ 
			$$\heva{&x=-7+3t\\&y=-10+4t \\&z=4-2t \\&z=0}.$$
			Giải hệ tìm được $N(-1; -2; 0)$. Lấy điểm $A(-7; -10; 4)\in d$. Gọi $H$ là hình chiếu vuông góc của $A$ lên mặt phẳng $(Oxy)$. Khi đó $H(-7; -10; 0)$.\\
			Đường thẳng $d'$ đi qua $N$ và $H$. Ta có $\overrightarrow{NH}=(-6; -8; 0)$. Một véc-tơ chỉ phương của $d$ là $(3; 4; 0)$.\\
			Phương trình đường thẳng $d'$ là $\heva{&x=-1+3t\\&y=-2+4t\\&z=0}$.
			\itemch Sai. Mặt phẳng $(P)$ cơ véc-tơ pháp tuyến $\overrightarrow{n}=(1;2;-3)$.\\
			Đường thẳng $d$ có véc-tơ chỉ phương $\overrightarrow{u}=(3;4;-2)$.\\
			Gọi $\Delta$ là đường thẳng cần lập phương trình. Theo giả thiết, ta có $[\overrightarrow{n},\overrightarrow{u}]=(8;-7;-2)$ là một véc-tơ chỉ phương của đường thẳng $\Delta$.\\
			Ta có $M(2; 2; -2)$ là giao của $d$ và $(P)$. Do đó $\Delta$ đi qua $M$.\\
			Vậy $\Delta$ có phương trình $\dfrac{x-2}{8}=\dfrac{y-2}{-7}=\dfrac{z+2}{-2}$.
		\end{itemchoice}
	}
\end{ex}

\Closesolutionfile{ans}
\indapan{3}{ans/ans-2-OTGK2-De1-DS}

\Opensolutionfile{ans}[ans/ans-2-OTGK2-De1-KQ]
\caukq
%Câu 1
\begin{ex}%[2D4N1-3]
	Biết hàm số $ F(x)$ là một nguyên hàm của hàm số $ f(x)=\sin x$ và thỏa mãn $F\left(\dfrac{\pi}{2}\right)=2$. Tính giá trị của $F(\pi)$.
	\shortans{$3$}
	\loigiai{
		Ta có $F(x)=\displaystyle\int\limits \sin x\mathrm{\,d}x=-\cos x+C$.\\
		$F\left(\dfrac{\pi}{2}\right)=2\Rightarrow C=2$.\\
		Khi đó $F(x)=-\cos x+2\Rightarrow F(\pi)=3$.
	}
\end{ex}
%Câu 2
\begin{ex}%[2D4H2-2]
	Biết $F(x)=x^{3}$ là một nguyên hàm của hàm số $f(x)$ trên $\mathbb{R}$. Tính giá trị của $\displaystyle\int\limits_{1}^{3}(1+f(x))\mathrm{\,d}x$.
	\shortans{$28$}
	\loigiai{
		Do $F(x)=x^3$ là một nguyên hàm của hàm số $f(x)$ trên $\mathbb{R}$ nên
		\[\displaystyle\int\limits_{1}^{3}(1+f(x))\mathrm{\,d}x=\displaystyle\int\limits_{1}^{3}\mathrm{\,d}x+\displaystyle\int\limits_{1}^{3}f(x)\mathrm{\,d}x=x\Big|_{1}^{3}+x^3\Big|_{1}^{3}=28.\]
	}
\end{ex}
%Câu 3
\begin{ex}%[2D4C3-3]
	\immini{
		Cho nửa lục giác đều $ABCD$ nội tiếp đường tròn đường kính $AD=8$ (tham khảo hình vẽ). Thể tích khối tròn xoay được tạo thành khi quay miền tứ giác $ABCD$ quanh đường thẳng $CD$ bằng $a\pi$. Tìm $a$.
	}
	{\begin{tikzpicture}[font=\footnotesize,line join=round, line cap=round, >=stealth,scale=0.55] 
			\def\r{4}
			\path ($(0,0)$) coordinate (O); 
			\foreach \x/\pos in {180/A, 120/B, 60/C, 0/D} \path ($(O)+(\x:\r)$) coordinate (\pos);
			\fill[pattern=north west lines,opacity=.6] (A)--(B)--(C)--(D)--(A);
			\draw (A)--(B)--(C)--(D)--(A) (O) circle (\r) ; 
			\fill (0,0) circle(1pt); 
			\foreach \x/\pos in {A/180, B/120, C/60, D/0} \fill (\x) circle(1pt) node[{shift=(\pos:0.25)}]{$\x$}; 
	\end{tikzpicture}}
	\shortans{$112$}
	\loigiai{
		\begin{center}
			\begin{tikzpicture}[font=\footnotesize,line join=round, line cap=round, >=stealth,scale=0.75] 
				\def\r{4}
				\def \xmin{-1}\def \xmax{7}\def \ymin{-1}\def \ymax{8} 
				\draw[->] (\xmin,0)--(\xmax,0) node[shift=(-90:0.25)] {$x$};
				\draw[->] (0,\ymin)--(0,\ymax) node[shift=(0:0.25)] {$y$};
				\fill (0,0) circle(1pt) node[shift=(-135:0.25)]{$O$}
				(4,0) circle(1pt) node[shift=(-45:0.2)]{$4$}
				(6,0) circle(1pt) node[shift=(-90:0.2)]{$6$}
				(0,{4*sqrt(3)}) circle(1pt) node[left]{$4\sqrt{3}$}
				(0,{2*sqrt(3)}) circle(1pt) node[left]{$2\sqrt{3}$};
				\path ($(0,0)$) coordinate (D)
				($(\r,0)$) coordinate (C)
				($(D)+(60:2*\r)$) coordinate (A)
				($(C)+(60:\r)$) coordinate (B);		
				\fill[pattern=north west lines,opacity=.6] (A)--(B)--(C)--(D)--(A);
				\draw[dashed] (A)--(C) (0,{2*sqrt(3)})--(B)--(6,0) (0,{4*sqrt(3)})--(A);
				\draw (D)--(A)--(B)--(C);
				\foreach \x/\pos in {A/90, B/0, C/-135, D/135} \fill (\x) circle(1pt) node[{shift=(\pos:0.25)}]{$\x$}; 
			\end{tikzpicture}
		\end{center}
		Đặt hệ trục tọa độ $Oxy$ như hình vẽ. Khi đó, ta có $A(4;4\sqrt{3})$, $B(6;2\sqrt{3})$, $C(4;0)$ và $D(0;0)$.\\
		Phương trình đường thẳng $AD$ là $y=\sqrt{3}x$.\\
		Phương trình đường thẳng $AB$ là $y=-\sqrt{3}x+8\sqrt{3}$.\\
		Phương trình đường thẳng $BC$ là $y=\sqrt{3}x-4\sqrt{3}$.\\
		Suy ra thể tích khối tròn xoay được tạo thành khi quay miền tứ giác $ABCD$ quanh đường thẳng $CD$ là
		\begin{align*}
			\displaystyle V&=\pi\int\limits_{0}^{4}\left(\sqrt{3}x\right)^2 \mathrm{\, d} x+ \pi\int\limits_{4}^{6}\left(\left(-\sqrt{3}x+8\sqrt{3}\right)^2-\left(\sqrt{3}x-4\sqrt{3}\right)^2\right)\mathrm{\, d} x\\ 
			&=\pi\int\limits_{0}^{4}3x^2 \mathrm{\, d} x+ 3\pi\int\limits_{4}^{6}\big(-8x+48\big)\mathrm{\, d} x\\ 
			&=\pi \cdot x^3\Big|_{0}^{4}+3\pi\big(-4x^2+48x\big)\Big|_{4}^{6}\\
			&=64\pi+48\pi=112\pi.
		\end{align*}
	}
\end{ex}
%Câu 4
\begin{ex}%[2H5H1-5]
	Trong không gian $Oxyz$, cho hai mặt phẳng $(\alpha)\colon x+2y+z-1$ và $(\beta)\colon 2x-y-z+2=0$. Gọi $E(1;b;0)$ với $b\in\mathbb{Z}$ và $E$ cách đều hai mặt phẳng $(\alpha)$ và $(\beta)$. Tìm $b$.
	\shortans{$-4$}
	\loigiai{
		Ta có
		\allowdisplaybreaks
		\begin{eqnarray*}
			\mathrm{d}\left(E, (\alpha)\right)=\mathrm{d}\left(E, (\beta)\right)&\Leftrightarrow& \dfrac{|2b|}{\sqrt{1^2+2^2+1^2}}=\dfrac{|4-b|}{\sqrt{2^2+(-1)^2+(-1)^2}}\\
			&\Leftrightarrow& |2b|=|4-b|\\
			&\Leftrightarrow& \hoac{&b=\dfrac{4}{3}(\text{loại})\\&b=-4.}
		\end{eqnarray*}
	}
\end{ex}
%Câu 5
\begin{ex}%[2H5N2-7]
	Trong không gian với hệ tọa độ $Oxyz$, cho phương trình đường thẳng $d\colon \dfrac{x-5}{1}=\dfrac{y+2}{1}=\dfrac{z-4}{\sqrt{2}}$ và phương trình mặt phẳng $(\alpha)\colon x-y+\sqrt{2}z-7 = 0$. Góc giữa đường thẳng $d$ và mặt phẳng $(\alpha)$ là bao nhiêu độ?
	\shortans{$30$}
	\loigiai{
		Đường thẳng $d$ có vec tơ chỉ phương là $\overrightarrow{u}=\left(1;1;\sqrt{2}\right)$.\\
		Mặt phẳng $(\alpha)$ có vec tơ pháp tuyến là $\overrightarrow{n}=\left(1;-1;-\sqrt{2}\right)$.\\
		Gọi $\varphi$ là góc giữa đường thẳng $d$ và mặt phẳng $(\alpha)$, ta có
		$$\sin \varphi =  \dfrac{\left|\overrightarrow{u}\cdot\overrightarrow{n}\right|}{\left|\overrightarrow{u}\right|\cdot\left|\overrightarrow{n}\right|}= \dfrac{|1\cdot 1 + 1 \cdot (-1) + \sqrt{2}\cdot \sqrt{2}|}{\sqrt{1^2+1^2+\sqrt{2}^2}\cdot\sqrt{1^2+(-1)^2+\sqrt{2}^2}}=\dfrac{1}{2} \Rightarrow \varphi = 30^\circ.$$
	}
\end{ex}
%Câu 6
\begin{ex}%[2H5C2-3]
	Trong không gian với hệ tọa độ $Oxyz$, một cabin cáp treo xuất phát từ điểm $A(10;3;0)$ và chuyển động đều theo đường cáp có véc-tơ chỉ phương là $\vv{u}=(2;-2;1)$ với tốc độ là $4{,}5$ m/s (đơn vị trên mỗi trục tọa độ là mét). Cabin dừng ở điểm $B$ có hoành độ $x_B=550$. Tìm độ dài quãng đường $AB$ (làm tròn kết quả đến hàng đơn vị của mét).
	\begin{center}
		\begin{tikzpicture}[line join=round,line cap=round,>=stealth,scale=0.75,font=\footnotesize]
			\tikzset{every node/.style={outer sep=-0.5mm}}
			\foreach \x/\y/\n in {0/0/O,5/0/y,0/4/z,-4/-4/x,-5/2.5/B} \coordinate (\n) at (\x,\y);
			\coordinate (C) at ($(O)!0.1!(y)$);
			\coordinate (D) at ($(O)!0.15!(x)$);
			\coordinate (A) at ($(C)+(D)-(O)$);
			\coordinate (E) at ($(O)!0.8!(x)$);
			\coordinate (M) at ($(A)!0.65!(B)$);
			\coordinate (v) at ($(A)!0.15!(B)$);
			\coordinate (u) at ($(O)+(v)-(A)$);
			\draw[->] (O)--(x) node[right] {\normalsize$x$};
			\draw[->] (O)--(y) node[below left] {\normalsize$y$};
			\draw[->] (O)--(z) node[below left] {\normalsize$z$};
			\draw[dashed,thin](C)--(A)--(D) (B)node[above]{$B$}--(E)node[below right]{$550$};
			\draw (E) node[left]{$x_B$};
			\draw (O) node[above right]{$O$};
			\draw (C) node[below]{$3$};
			\draw (D) node[left]{$10$};
			\fill (A) circle (1.5pt) node[below right] {$A(10;3;0)$};
			\draw[very thick] (A)--(B);
			\draw[->,very thick,red] (O)--(u)node[above]{$\vv{u}$};
			\fill[red] (M) circle (2pt) node[above right] {$M$};
			\draw (M) .. controls ($(M)+(-0.2,-0.1)$) .. ($(M)+(0,-0.2)$) node[below,yshift=1mm]{{\large \faCube}};
			%			\draw ($(O)+(0,-4)$) node[below]{\textit{Hình 35}};
		\end{tikzpicture}
	\end{center}
	\shortans{$810$}
	\loigiai{
		\begin{itemize}
			\item Phương trình chính tắc của đường cáp là $\dfrac{x-10}{2}=\dfrac{y-3}{-2}=\dfrac{z}{1}$.
			\item Do tốc độ chuyển động của cabin là $4{,}5$ m/s nên độ dài $AM=4{,}5t$ (m). \\
			Vì vậy $\left|\vv{AM}\right|=4{,}5t$ ($t\ge0$).\\
			Do hai véc-tơ $\vv{AM}$ và $\vv{u}$ là cùng hướng nên $\vv{AM}=k\vv{u}$ với $k\in\mathbb{R}^+$.\\
			Suy ra $\left|\vv{AM}\right|=k\left|\vv{u}\right|=k\cdot\sqrt{2^2+(-2)^2+1^2}=3k$. Do đó $3k=4{,}5t\Leftrightarrow k=\dfrac{3t}{2}$. \\
			Vì thế, ta có $\vv{AM}=\dfrac{3t}{2}\vv{u}=\left(3t;-3t;\dfrac{3t}{2}\right)$.\\
			Gọi tọa độ của điểm $M$ là $(x_M;y_M;z_M)$.\\
			Do $\vv{AM}=(x_M-x_A;y_M-y_A;z_M-z_A)=\left(3t;-3t;\dfrac{3t}{2}\right)$ nên $$\heva{&x_M=3t+x_A\\&y_M=-3t+y_A\\&z_M=\dfrac{3t}{2}+z_A}\Leftrightarrow\heva{&x_M=3t+10\\&y_M=-3t+3\\&z_M=\dfrac{3t}{2}.}$$
			Vậy điểm $M$ có tọa độ là $\left(3t+10;-3t+3;\dfrac{3t}{2}\right)$.
			\item Do $x_B=550$ nên $3t+10=550$, tức là $t=180$ (s). Do đó, ta có điểm $B(550;-537;270)$.\\
			Vậy $AB=\sqrt{(550-10)^2+(-537-3)^2+(270-0)^2}=\sqrt{656100}=810$ (m).
		\end{itemize}
	}
\end{ex}
\Closesolutionfile{ans}
\indapan{6}{ans/ans-2-OTGK2-De1-KQ}